%%%%%%%%%%%%%%%%%%%%%%%%%%%%%%%%%%%%%%%%%%%%%%%%%%%%%
%
% This is the template for the EGAS 55 conference abstracts.
% Please, insert your name and affiliations, text of the abstract and references
% at the locations in the file where they are mentioned.
% Do not change the part defining the format of the file.
%%%%%%%%%%%%%%%%%%%%%%%%%%%%%%%%%%%%%%%%%%%%%%%%%%%%%
%
%%%%%%%%%%%%%%%%%%%%   Format, do not change  %%%%%%%%%%%%%%%%%
%

% LaTeX file to be processed by pdflatex

\documentclass[11pt,a4paper]{article}

\usepackage{times}
\usepackage{graphicx}

%\usepackage{draftwatermark}


\setlength{\hoffset}{0.0cm}
\setlength{\voffset}{0.0cm}

\setlength{\oddsidemargin}{-0.8cm}
\setlength{\marginparwidth}{-0.5cm}

\setlength{\headsep}{0.0cm}
\setlength{\headheight}{0.0cm}
\setlength{\topskip}{0.4cm}
\setlength{\topmargin}{-0.5cm}
\setlength{\footskip}{1.0cm}

\setlength{\textwidth}{16.9cm}
\setlength{\textheight}{24.7cm}

\setlength{\parindent}{0em}


\newcommand{\abstracttitle}[1]{\parbox{\textwidth}
                              {\setlength{\baselineskip}{19pt}%
                               \begin{center}{\large\bfseries #1}\end{center}}}
\newcommand{\authors}[1]{\begin{center}#1\end{center}}
\newcommand{\addresses}[1]{\begin{center}#1\end{center}}

\renewcommand{\labelenumi}{[\arabic{enumi}]}

\pagestyle{empty}
% \pagestyle{headings}
\markright{EGAS 55}

%%%%%%%%%%%%%%   End format   %%%%%%%%%%%%%%%%%


%%%%%%%%%%%%%%%%%%%%%%%%%%%%%%%%%%%%%%%%%%%
%             Editable part
%%%%%%%%%%%%%%%%%%%%%%%%%%%%%%%%%%%%%%%%%%%
%

\begin{document}
%
% Title of the abstract
\abstracttitle{USOC: Ultra-precise, shock-resistant optical clocks}
%
% Authors and affiliation
\authors{\textit{\underline{Thomas Walker*}, Ludovico Iannizzotto Venezze, Leonie Hawkins, Charles Baynham, Richard Hobson, Oliver Buchmueller}}

\addresses{{\small London, England\\ \vspace{0.2cm}
(*) t.walker@imperial.ac.uk}}

\setlength{\parskip}{\bigskipamount}
%%%%%%%%%%%%%%%%%%%%%%%%%%%%%%%
% Text of your abstract
%%%%%%%%%%%%%%%%%%%%%%%%%%%%%%%

Optical clocks provide unparalelled precision for the measurement of time, and are set to replace the definition of the second. 
However, many potenial applications of these devices, such as for global positioning or in commication networks, require deployment outside the lab in noisy environments. 
In particular, mechanical shock can greatly impact short-term stability. This can be mitigated by using a continuous measurement scheme, rather than the conventional pulsed schemes.
By continuously loading cold atoms into a measurement region, measurement dead-time can be eliminated and so the associated loss of precision can be reduced. 


Important remarks:
\begin{itemize}
\item Abstracts should be written in English.

\item The abstract should fit on one A4 page (210 mm $\times$ 297 mm) and should come without a page number.

\item  It should be submitted in PDF format  through the conference website, using this template. 

\item Figures could be also included.

\item   Files larger than 5 MB will be rejected.

\end{itemize}



 \vspace{0.3cm}

% Uncomment the following if you wish to include a figure in the abstract
%
%\begin{figure}[h]
%\begin{center}
%\includegraphics*[height=2.0cm]{Figure.png}
%\end{center}
%\caption{The caption for the figure should go here.}
%\end{figure}

\setlength{\parskip}{\smallskipamount}
\noindent 

% Your references 
\begin{thebibliography}{10}
\bibitem{Quixote_1605} 
\textit{El ingenioso hidalgo don Quijote de la Mancha}, M. de Cervantes Saavedra, (First Edition, 1605).

\bibitem{grossman} 
\textit{Don Quixote}, M. de Cervantes (Author), E. Grossman (Translator), (HarperCollins Publishers Inc., New York, 2003) .

\bibitem{Sancho} 
S. Panza, and M. Cervantes, Phys. Rev. X {\bf 0000}, 1 (1605).


\end{thebibliography}
% End your references

\end{document}
